\documentclass{article}
\usepackage[utf8]{inputenc}
\usepackage{listings}
\usepackage{xcolor}
\usepackage{hyperref}
\usepackage{geometry}
\geometry{a4paper, margin=1in}

\title{Detailed Report: Code Analysis of \texttt{4\_task.R}}
\author{}
\date{\today}

\definecolor{codegreen}{rgb}{0,0.6,0}
\definecolor{codegray}{rgb}{0.5,0.5,0.5}
\definecolor{codepurple}{rgb}{0.58,0,0.82}
\definecolor{backcolour}{rgb}{0.95,0.95,0.92}

\lstdefinestyle{mystyle}{
    backgroundcolor=\color{backcolour},   
    commentstyle=\color{codegreen},
    keywordstyle=\color{magenta},
    numberstyle=\tiny\color{codegray},
    stringstyle=\color{codepurple},
    basicstyle=\ttfamily\footnotesize,
    breakatwhitespace=false,         
    breaklines=true,                 
    captionpos=b,                    
    keepspaces=true,                 
    numbers=left,                    
    numbersep=5pt,                  
    showspaces=false,                
    showstringspaces=false,
    showtabs=false,                  
    tabsize=2
}
\lstset{style=mystyle}

\begin{document}

\maketitle

\section{Introduction}
This report provides a detailed, line-by-line analysis of the \texttt{4\_task.R} file, which corresponds to Task 4 of the project. The primary goal of this script is to train a Multinomial Logistic Regression model to predict the probability of a data point belonging to a specific cluster (identified in Task 3) based solely on the hour of the day. To correctly model the circular nature of time, Fourier transformations (harmonics) are utilized.

\section{Line-by-Line Code Analysis}

\subsection{Data Preparation and Feature Engineering}
\begin{lstlisting}[language=R]
# ==============================================================================
# TASK 4: Advanced Classification Model (Multinomial Logistic Regression)
# ==============================================================================

cat("\nExecuting Task 4: Fitting multinomial model with Fourier-transformed hours...\n")

# Assicuriamoci che la colonna map_cluster sia un fattore
df$map_cluster <- as.factor(df$map_cluster)

# 1. Feature Engineering: Trasformata di Fourier per l'ora (circolare)
# Aggiungiamo l'armonica principale (24h) e la secondaria (12h) per catturare eventuali doppi picchi giornalieri
df$sin24 <- sin(2 * pi * df$hour / 24)
df$cos24 <- cos(2 * pi * df$hour / 24)
df$sin12 <- sin(2 * pi * 2 * df$hour / 24)
df$cos12 <- cos(2 * pi * 2 * df$hour / 24)
\end{lstlisting}
\begin{itemize}
  \item \textbf{Lines 1-8}: Introductory comments and ensuring the \texttt{map\_cluster} column (the target variable) is strictly a \texttt{factor}. This categorical typing is required by the multinomial regression function.
  \item \textbf{Lines 10-15}: This is a crucial feature engineering step. Since hours are cyclical (hour 23 is adjacent to hour 0), using the raw numerical value of the hour in a linear model is incorrect. The time is transformed into a continuous circular representation using Fourier terms (sine and cosine waves).
        \begin{itemize}
          \item \texttt{sin24} and \texttt{cos24}: The fundamental harmonic representing a full 24-hour cycle.
          \item \texttt{sin12} and \texttt{cos12}: The second harmonic representing a 12-hour cycle. Including this allows the model to capture more complex daily patterns, such as a double-peak behavior (e.g., morning and evening consumption peaks).
        \end{itemize}
\end{itemize}

\subsection{Model Training}
\begin{lstlisting}[language=R, firstnumber=17]
# 2. Addestramento del modello logistico multinomiale
cat("Fitting multinomial logistic regression...\n")
# Usiamo la funzione multinom del pacchetto nnet (caricato nel setup)
# trace = FALSE per nascondere l'output di ottimizzazione a schermo
multinom_model <- multinom(map_cluster ~ sin24 + cos24 + sin12 + cos12, data = df, trace = FALSE)
\end{lstlisting}
\begin{itemize}
  \item \textbf{Line 21}: The Multinomial Logistic Regression model is trained using the \texttt{multinom()} function from the \texttt{nnet} package. The formula indicates that \texttt{map\_cluster} is being predicted using the four newly created Fourier features. The \texttt{trace = FALSE} argument suppresses the console output generated by the internal optimization algorithm.
\end{itemize}

\subsection{Continuous Grid Prediction}
\begin{lstlisting}[language=R, firstnumber=23]
# 3. Previsione su una griglia continua di ore
cat("Predicting probabilities over a continuous 24h grid...\n")
# Creiamo 200 punti per avere una curva molto morbida
hour_grid <- seq(0, 23, length.out = 200)

grid_df <- data.frame(
  hour = hour_grid,
  sin24 = sin(2 * pi * hour_grid / 24),
  cos24 = cos(2 * pi * hour_grid / 24),
  sin12 = sin(2 * pi * 2 * hour_grid / 24),
  cos12 = cos(2 * pi * 2 * hour_grid / 24)
)

pred_probs <- predict(multinom_model, newdata = grid_df, type = "probs")
\end{lstlisting}
\begin{itemize}
  \item \textbf{Lines 25-26}: A continuous grid of hours (\texttt{hour\_grid}) is created with 200 evenly spaced points from 0 to 23. This high resolution allows us to plot a smooth probability curve later.
  \item \textbf{Lines 28-34}: A new data frame, \texttt{grid\_df}, is created. It computes the identical Fourier transformations (sine and cosine for 24h and 12h periods) for the continuous \texttt{hour\_grid} points.
  \item \textbf{Line 37}: The trained \texttt{multinom\_model} is used to predict the probabilities (\texttt{type = "probs"}) over the new \texttt{grid\_df}. This returns a matrix \texttt{pred\_probs} where each row is a time point from the grid, and each column corresponds to a cluster's probability (the rows sum to 1).
\end{itemize}

\subsection{Data Formatting for Plotting}
\begin{lstlisting}[language=R, firstnumber=39]
# Convertiamo in dataframe e formattiamo i dati in formato 'long' per ggplot
prob_df <- as.data.frame(pred_probs)

if (best_model == 2) {
  prob_df <- data.frame(`1` = 1 - pred_probs, `2` = pred_probs)
}
colnames(prob_df) <- paste0("Cluster_", 1:ncol(prob_df))
prob_df$hour <- hour_grid

# Usa pivot_longer da tidyr (caricato nel setup) per preparare i dati per ggplot
library(tidyr)
prob_long <- pivot_longer(prob_df, cols = starts_with("Cluster"), names_to = "Cluster", values_to = "Probability")
# Pulisce i nomi per corrispondere visivamente ai cluster di df (es. "1", "2")
prob_long$Cluster <- sub("Cluster_", "", prob_long$Cluster)
prob_long$Cluster <- as.factor(prob_long$Cluster)
\end{lstlisting}
\begin{itemize}
  \item \textbf{Lines 40-48}: The prediction matrix is converted into a data frame \texttt{prob\_df}. An edge case is handled at lines 43-45: if the optimal model selected in Task 1 had exactly $H=2$ components, the \texttt{predict} function for multinomial regression returns a single vector (the probability of the second class) instead of a matrix. In this scenario, the code manually reconstitutes the two columns. Columns are then renamed systematically, and the continuous hour variable is appended.
  \item \textbf{Lines 50-54}: The \texttt{tidyr} package is used to pivot the data frame from wide to long format (\texttt{prob\_long}), which is the required structure for \texttt{ggplot2}. The "Cluster\_" prefix is removed from the names for a cleaner legend display, and the column is converted to a factor.
\end{itemize}

\subsection{Generating and Saving the Final Plot}
\begin{lstlisting}[language=R, firstnumber=56]
# 4. Generazione del grafico
cat("Generating probability plot...\n")
p_multinom <- ggplot(prob_long, aes(x = hour, y = Probability, fill = Cluster)) +
  geom_area(alpha = 0.85, color = "white", linewidth = 0.2) +
  theme_minimal() +
  scale_x_continuous(breaks = seq(0, 23, by = 2)) +
  labs(
    title = "Multinomial Regression: Predicted Probability of Clusters",
    subtitle = "Continuous model based on Fourier-transformed hours",
    x = "Hour of the Day (0-23)",
    y = "Predicted Probability",
    fill = "Cluster"
  ) +
  theme(
    plot.title = element_text(face = "bold", size = 14),
    axis.title = element_text(face = "bold"),
    legend.position = "bottom"
  )

print(p_multinom)

# Salvataggio del grafico
ggsave("images/multinomial_probabilities.png", plot = p_multinom, width = 8, height = 5, dpi = 300)
cat("\nPlot salvato con successo come 'images/multinomial_probabilities.png' nella cartella del progetto.\n")
\end{lstlisting}
\begin{itemize}
  \item \textbf{Lines 61-78}: A stacked area chart (\texttt{geom\_area}) is constructed using \texttt{ggplot2}. Because the predicted probabilities sum to exactly 1 at every given point in time, a stacked area plot is an ideal choice. It provides a smooth, continuous version of the stacked bar chart seen in Task 3, allowing for intuitive visual comparison. The plot features customized titles, a minimal theme, and adjusted x-axis breaks.
  \item \textbf{Lines 80-82}: The plot is printed and exported as a high-resolution PNG image at \texttt{images/multinomial\_probabilities.png}.
\end{itemize}

\section{Conclusion}
Task 4 successfully elevates the analysis from basic descriptive clustering to predictive modeling. By properly treating time as a circular variable using Fourier harmonics, the Multinomial Logistic Regression model effectively smooths and interpolates the cluster assignment probabilities over the entire day. The final visualization serves as an elegant summary of how the likelihood of observing different load factor behaviors evolves continuously across the 24-hour cycle.

\end{document}
